\chapter{Numerical methods for driven problems}\label{ch:drivenmethods}

\section{From a drive to harmonics}\label{sec:sampling}

\code{PeriodicHamiltonian.from_parameters} evaluates the model constructor at
$N$ equally spaced times $t_j=jT/N$ and Fourier transforms every operator coefficient,
\begin{equation}\label{eq:fft}
  H_k=\frac1N\sum_{j=0}^{N-1}H(t_j)\,\ee^{-2\pi\ii jk/N},\qquad 0\le k\le\lceil N/2\rceil-1 .
\end{equation}
This is exact for Hamiltonians with no harmonics at or above $N/2$, including a
monochromatic modulation of a parameter on which the Hamiltonian depends linearly. For the phase drive the Bessel functions decay
faster than exponentially once $k>A_\phi/4$, so a modest $N$ suffices. The construction is
checked at the midpoints $t_{j+1/2}$, where the model is rebuilt and compared with the
Fourier series. The largest $\ell_1$ norm of the coefficient differences bounds the
operator-norm error of the interpolated Hamiltonian and must stay below
\code{sampling_tolerance}. The observables of the model, e.g.\ $\partial H/\partial\phi$,
are sampled in the same way, since they too depend on time under a phase drive.

\section{Sambe-space diagonalization}\label{sec:sambe-method}

\code{solve_floquet} truncates $Q$ in \cref{eq:sambe-matrix} to $\abs m\le M$ (option
\code{floquet_cutoff}) and diagonalizes it. The many-body space in each block is either the full
QP Fock space of the sector (\code{basis="fock"}) or the static basis of the $K$ lowest
eigenstates of the undriven Hamiltonian (\code{basis="eigen"}, $K$ set by
\code{basis_states}). In the static basis the harmonics become the $K\times K$ matrices
$\matel a{H_k}b$ between the static eigenstates. The Sambe dimension is $(2M+1)d$, where
$d$ is the dimension of the Fock space of the sector or $d=K$. Small problems are
diagonalized densely. Large Fock-space problems use shift-invert Arnoldi iterations
(ARPACK)~\cite{lehoucq1998arpack}, with a sparse LU factorization of $Q-\sigma$ for each
target quasienergy $\sigma$~\cite{li2005superlu}; by default the targets are the static
energies of the tracked states.

\paragraph{A rigorous error bound.}
Let $v$ be a normalized eigenvector of the truncated matrix $Q_M$ with eigenvalue $\lambda$, and
embed it in the infinite Sambe space. The residual of the embedded vector has two parts:
the window residual $r=\norm{(Q_M-\lambda)v}$, which measures the eigensolver accuracy, and
the \emph{leakage} $\kappa=\norm{\sum_{m'}H_{m-m'}v_{m'}}$ into the omitted blocks
$\abs m>M$. Because the infinite Sambe operator is self-adjoint, it has an eigenvalue within
\begin{equation}\label{eq:residual-bound}
  \abs{\qe-\lambda}\le\sqrt{r^2+\kappa^2}
\end{equation}
of $\lambda$ (\cref{app:residual}). This combined residual is reported as
\code{residuals}; the window residual must stay below \code{residual_tolerance}. The bound is
of first order: for an isolated quasienergy the actual error is typically of second order in
$\kappa$, and often much smaller than the bound. In the static basis, \code{fock_residuals}
evaluates the same residual for the Fock-space Sambe operator and thus tests the
convergence in $K$.

\paragraph{Branches and tracking.}
Each Floquet state appears in many replicas (\cref{sec:sambe}). The solver keeps one replica
per state, the one whose harmonic centroid $\expval m$ is smallest in magnitude. It
recognizes the other replicas by their quasienergy, shifted by a multiple of $\Omega$, and
by their overlap after a shift of the harmonic index; for degenerate multiplets it projects
onto all retained states with that quasienergy. With \code{selection="tracked"} (the default), the
lowest \code{states} static eigenstates are assigned to Floquet states by maximizing the
total static weight $s_{\alpha a}=\sum_m\abs{\braket a{u_{\alpha,m}}}^2$; each assigned weight must
exceed \code{tracking_threshold}. Near a resonance the eigenpairs closest to one shift-invert
target can be replicas of the state assigned at another target. The solver therefore doubles
the number of eigenpairs requested at that target until the corresponding static state is assigned, up to
\code{max_eigenpairs_per_target}. With
\code{selection="all"}, every Floquet state in the window is returned. Complete sets, needed
for the sum rule of \cref{ch:probes}, require the dense solver, because shift-invert finds
only the states near explicit \code{targets}. The reported quasienergy is placed on a
definite branch:
\begin{itemize}
\item \code{branch="label"} (default for tracked states) chooses the replica in which the
component on the assigned static state $a$ has its harmonic centroid in $[-\tfrac12,\tfrac12)$,
so that $\qe_\alpha\to E_a$ as the drive vanishes and $\qe_\alpha$ varies continuously
through a one-photon resonance;
\item \code{branch="centroid"} (default for \code{selection="all"}) uses the replica with
total centroid in $[-\tfrac12,\tfrac12)$;
\item \code{folded} gives every quasienergy in the zone $[c-\Omega/2,c+\Omega/2)$,
$c=$\code{zone_center}.
\end{itemize}

\begin{caution}
At a resonance the dressed states are equal mixtures of two static states, and the
labels may swap from one side of the resonance to the other. Plot both branches together,
as in the rotating frame of \cref{fig:tls-weak}, to follow the complete resonant pair.
\end{caution}

\section{The one-period propagator}\label{sec:propagator}

\code{floquet_propagator} integrates the Schr\"odinger equation over one period with
$H(t)=\sum_kH_k\ee^{\ii k\Omega t}$, without a Floquet cutoff, and diagonalizes $U(T,0)$. The
integrator is the fourth-order commutator-free Magnus
scheme~\cite{blanes2006fourth,alvermann2011high,blanes2009magnus} (\cref{app:cfm4}),
\begin{equation}\label{eq:cfm4}
  U(t+h,t)=\ee^{-\ii h(\alpha_1H_1+\alpha_2H_2)}\,\ee^{-\ii h(\alpha_2H_1+\alpha_1H_2)},\qquad
  \alpha_{1,2}=\frac{3\mp2\sqrt3}{12},\quad H_{1,2}=H\bigl(t+(\tfrac12\mp\tfrac{\sqrt3}6)h\bigr),
\end{equation}
where $h$ is the time step. The exponentials are evaluated by eigendecomposition, so that
each step is exactly unitary. The eigenphases of $U(T,0)$ give the folded quasienergies. The
mean energies \eqref{eq:mean-energy}, accumulated along a second propagation, select the
central replica, and the Fourier transform of $\ket{u(t)}$ on the time grid gives the
harmonics. With \code{step_check=True} the calculation is repeated with half the time step,
and $\abs{\qe_h-\qe_{h/2}}/15$ estimates the remaining error of the time stepping. The
propagator approach is also used by QuTiP's Floquet
module~\cite{johansson2012qutip,lambert2026qutip}. In \pkg{} it serves as an independent
check of convergence with the Floquet cutoff and as an alternative method in a static basis.

\section{The high-frequency expansion}\label{sec:van-vleck}

For $\Omega$ large compared with the many-body spectral width, the Sambe blocks with
$m\ne0$ are far from the $m=0$ block and can be eliminated perturbatively. The result is the
static effective Hamiltonian $\Heff$ introduced in \cref{ch:floquet}, whose eigenvalues approximate the
quasienergies. To second order in $1/\Omega$, the van Vleck form of this
expansion~\cite{eckardt2015high,bukov2015universal,goldman2014periodically,mikami2016brillouin}
reads (\cref{app:vanvleck})
\begin{equation}\label{eq:van-vleck}
  \Heff=H_0+\frac1\Omega\sum_{m\ge1}\frac{[H_m,H_{-m}]}m
  +\frac1{\Omega^2}\Bigl(\sum_{m\ne0}\frac{[[H_{-m},H_0],H_m]}{2m^2}
  +\sum_{m\ne0}\sum_{m'\ne0,m}\frac{[[H_{-m'},H_{m'-m}],H_m]}{3mm'}\Bigr).
\end{equation}
The first-order term vanishes when every $H_k$ commutes with $H_{-k}=H_k\dg$. This holds
for any drive that becomes even in time after a shift of the time origin, such as a
single-parameter cosine with any phase $\theta$, for which $H_k=\ee^{\ii k\theta}g_k$ with
Hermitian $g_k$. A finite first-order term requires, for example, several parameters driven
with relative phases. The commutators are evaluated with the full
fermionic operators, before any QP truncation. The result is an ordinary
\code{Hamiltonian} in the same coordinates (\code{periodic.van_vleck(order)}), which both the
QP and the DMRG solver can solve. The error at order $n$ scales as $\Omega^{-(n+1)}$
(\cref{fig:highfreq}). Because the bath bandwidth enters the many-body spectral width, the
expansion serves as a check for small models and as a conceptual tool in the high-frequency
regime. Treat subgap resonances at $\Omega\sim\Delta$ with Sambe diagonalization or the propagator.

\section{The rotating-wave model}\label{sec:rwa}

Near a resonance only a few blocks of $Q$ matter. \code{rwa_model(periodic, states, photons)}
restricts $Q$ to the pairs $\ket{a_i,n_i}$ of a static state $a_i$ and a photon number $n_i$,
\begin{equation}\label{eq:rwa}
  Q^{\mathrm{RWA}}_{ij}=\bigl(H_{n_i-n_j}\bigr)_{a_ia_j}+n_i\Omega\,\delta_{ij},
\end{equation}
with matrix elements in the static eigenbasis. The photon numbers are chosen so that the
diagonal entries $E_{a_i}+n_i\Omega$ are nearly degenerate. For example, absorption from $a$
to $b$ with $E_b-E_a\approx\Omega$ uses \code{states=[a, b]} and \code{photons=[0, -1]}. At
resonance the splitting is $2\abs{(H_{n_a-n_b})_{ab}}$. The model neglects the
counter-rotating terms (Bloch--Siegert shift) and the AC Stark shifts from all other
states, both of second order in the drive. It is exact for a two-level system with a
circularly polarized drive.

\section{Choosing a method}

\begin{table}[htb]
\centering\small
\caption{Methods for Floquet quasienergies in \pkg{}.}\label{tab:methods}
\begin{tabularx}{\textwidth}{>{\raggedright\arraybackslash}p{3.1cm}>{\raggedright\arraybackslash}p{3.6cm}X}
\toprule
Method & Controlled by & Use for \\
\midrule
Sambe, Fock basis & $M$, QP cutoff, bath & reference results; moderate Fock dimensions; shift-invert for tracked states \\
Sambe, static basis & $M$, $K$ & large Fock spaces, scans over $\Omega$ or amplitude; check \code{fock_residuals} \\
Propagator & time steps, basis & independent check without Floquet cutoff; strong slow drives \\
Van Vleck & order in $1/\Omega$ & $\Omega$ above the many-body width; effective Hamiltonians, DMRG \\
RWA & chosen states & interpretation of resonances, Rabi splittings \\
\bottomrule
\end{tabularx}
\end{table}

Convergence must be checked separately with respect to the Floquet cutoff $M$, the size
$K$ of the static basis, the QP cutoff, the bath, the number of drive samples and, for the
propagator, the number of time steps. \Cref{ch:ex-numerics} works through each of them.

\section{Validation}\label{sec:validation-driven}

The driven solvers are checked against independent results:
\begin{itemize}
\item propagation of the exact $H(t)$ of a phase- and gate-driven junction in the physical
electron basis, with a different Magnus integrator and without Fourier series; it agrees with
the QP propagator and the Fock-space Sambe results to $10^{-8}$;
\item the exactly solvable circularly driven two-level system, for Sambe, propagator, RWA and
the analytic van Vleck operators;
\item Bessel-function harmonic weights of a gate-driven isolated dot
(\cref{ch:ex-dot}) and Jacobi--Anger harmonics of the phase drive;
\item the gauge identity \eqref{eq:phase-to-level} between a phase drive and a level drive;
\item the $\Omega^{-(n+1)}$ error scaling of the van Vleck expansion;
\item the Rabi splitting at an $S$--$S^*$ resonance against the RWA and the static matrix
element.
\end{itemize}
